Die automatisierte Nutzung von Produktinformationen setzt voraus, dass inhaltlich heterogene und unterschiedlich dargestellte Angaben in einer einheitlichen, maschinell verarbeitbaren Form vorliegen. Bei Produktseiten im Web besteht jedoch eine Lücke zwischen der für Menschen verständlichen Darstellung und den strukturellen Anforderungen nachgelagerter Informationssysteme. Die Überführung solcher Inhalte erfordert daher nicht nur eine technische Formatänderung, sondern auch die Erkennung, fachliche Zuordnung und Prüfung der relevanten Informationen. Große Sprachmodelle eröffnen hierfür neue Möglichkeiten, deren Eignung und Grenzen unter kontrollierten Bedingungen untersucht werden müssen.

Dieses Kapitel führt in den Untersuchungsgegenstand der Arbeit ein und grenzt ihn gegenüber angrenzenden Aufgaben ab. Abschnitt 2.1 leitet aus der allgemeinen Bedeutung standardisierter Repräsentationen die konkrete Problemstellung der Transformation heterogener Lebensmittel-Produktdaten ab und begründet den Untersuchungsbedarf. Darauf aufbauend formuliert Abschnitt 2.2 die Zielsetzung, die Forschungsfragen und den betrachteten Untersuchungsumfang. Abschnitt 2.3 erläutert abschließend den Aufbau der Arbeit und ordnet die nachfolgenden Kapitel in die Argumentations- und Untersuchungslogik ein.

\subsection{Problemstellung und Motivation}

Bis in die Mitte des 20. Jahrhunderts war der Transport von Stückgut durch eine Vielzahl unterschiedlicher Behältnisse geprägt. Waren erreichten die Häfen in Säcken, Kisten, Fässern oder einzelnen Kartons und mussten beim Wechsel zwischen Lastkraftwagen, Lagerhaus, Schiff und Bahn wiederholt sortiert, verladen und gesichert werden. Der ab 1956 erprobte Frachtcontainer veränderte dieses Prinzip nicht dadurch, dass die transportierten Güter vereinheitlicht wurden. Stattdessen erhielten unterschiedlichste Waren eine standardisierbare äußere Form, die zwischen mehreren Verkehrsträgern übergeben werden konnte. Der Inhalt blieb verschieden, während die gemeinsame Form seine Handhabung an den Übergabepunkten vereinfachte \cite{levinson2016box}.

Ein vergleichbares Grundproblem tritt bei der Verarbeitung digitaler Informationen auf. Auch hier müssen die zugrunde liegenden Gegenstände nicht gleichartig sein. Informationssysteme benötigen jedoch eine gemeinsame Repräsentation, um Daten speichern, vergleichen und an nachgelagerte Prozesse übergeben zu können. Anders als beim physischen Transport genügt es dabei nicht, vorhandene Inhalte unverändert in einen neuen Behälter zu legen. Die relevanten Informationen müssen zunächst erkannt, fachlich interpretiert und den vorgesehenen Feldern einer Zielstruktur zugeordnet werden.

Diese Herausforderung zeigt sich insbesondere bei Produktinformationen aus dem Web. Webseiten sind durch HTML technisch gegliedert, ihre fachlichen Inhalte folgen jedoch keinem über alle Anbieter und Produkte einheitlichen Schema. Produktbezeichnungen, Marken, Zutatenlisten, Nährwerte, Produktgruppen oder Verpackungsgrößen können in Überschriften, Tabellen, Listen, frei formulierten Texten oder eingebetteten Metadaten auftreten. Vergleichbare Informationen werden unterschiedlich bezeichnet, an verschiedenen Positionen dargestellt oder in abweichender Granularität angegeben. Gleichzeitig enthalten Produktseiten Navigationsbereiche, Shopinformationen, Empfehlungen, Preisangaben und weitere Inhalte, die nicht oder nur mittelbar zum betrachteten Produkt gehören. Die Überführung solcher Daten in eine gemeinsame Repräsentation ist damit eine Aufgabe der Datenintegration und Informationsextraktion \cite{doan2012principles,ferrara2014web}.

Im Projektkontext der vorliegenden Arbeit werden Lebensmittel-Produktseiten bereits durch bestehende Crawling-Prozesse erfasst. Die entstehenden Rohdatenobjekte enthalten unter anderem den HTML-Inhalt einer Webseite sowie technische Metadaten. Sie stellen die vorhandenen Produktinformationen bereit, bilden diese jedoch nicht unmittelbar in der Struktur ab, die das nachgelagerte Informationssystem benötigt. Dem heterogenen Ausgangsmaterial steht ein fest definiertes, typisiertes und verschachteltes Zielschema gegenüber. Zwischen beiden Seiten besteht daher eine Transformationslücke: Relevante Informationen müssen aus den Quelldaten extrahiert, dem passenden Zielattribut zugeordnet, in die verlangte Darstellung überführt und vor der weiteren Verwendung geprüft werden.

Eine manuelle Bearbeitung kann diese Zuordnung für einzelne Datensätze leisten, verursacht bei wachsenden Datenbeständen jedoch einen entsprechend steigenden personellen Aufwand. Fehlerhafte oder unvollständige Transformationen wirken sich zudem auf die Nutzbarkeit der erzeugten Daten aus. Werden beispielsweise eine Marke und ein Händler verwechselt, eine Verpackungsgröße aus einer Grundpreisangabe abgeleitet oder nicht vorhandene Informationen ergänzt, kann das Ergebnis formal strukturiert und dennoch fachlich ungeeignet sein. Für nachgelagerte Aufgaben wie Suche, Vergleich, Analyse oder Systemintegration sind deshalb sowohl eine konsistente Repräsentation als auch korrekte und hinreichend vollständige Inhalte erforderlich \cite{wang1996beyond,batini2016data}.

Ein etablierter Lösungsweg besteht in regelbasierten Extraktionsverfahren. Sie verwenden explizite Vorgaben zu erwarteten Dokumentstrukturen, Zeichenmustern, Schlüsselwörtern und fachlichen Zusammenhängen. Bei unverändertem Regelwerk ermöglichen sie grundsätzlich reproduzierbare und nachvollziehbare Entscheidungen. Ihre Abdeckung hängt jedoch davon ab, welche Darstellungsformen, Schreibweisen und Sonderfälle im Regelwerk berücksichtigt wurden. Neue Quellenstrukturen oder bislang unbekannte Formulierungen können deshalb Ergänzungen und Anpassungen erforderlich machen \cite{ferrara2014web}. Bei einem umfangreichen Zielschema mit unterschiedlichen Attributtypen steigt zugleich die Anzahl der fachlichen Regeln, die entwickelt, abgestimmt und gepflegt werden müssen.

Große Sprachmodelle stellen einen alternativen Ansatz für diese Transformationsaufgabe dar. Sie können sprachliche und strukturelle Hinweise innerhalb eines gemeinsamen Eingabekontexts verarbeiten und durch Anweisungen auf eine vorgegebene Zielstruktur ausgerichtet werden. Dadurch erscheint es möglich, heterogene Produktinformationen zu interpretieren, ohne jede Darstellungsvariante als separate technische Regel abzubilden. Die Ausgabe eines strukturierten Objekts belegt jedoch noch nicht, dass dessen Inhalt korrekt, vollständig oder mit dem vorgesehenen Schema vereinbar ist. Untersuchungen zur generativen Informationsextraktion zeigen, dass die erzielte Qualität unter anderem vom betrachteten Attribut und von der konkreten Aufgabenbeschreibung abhängen kann \cite{data_extraction_ai}. Eine unabhängige Prüfung der syntaktischen Verarbeitbarkeit, der strukturellen Konformität und der fachlichen Werte bleibt daher erforderlich.

Aus dieser Ausgangslage ergibt sich der Bedarf an einem Lösungsansatz, der die flexible Verarbeitung heterogener Produktinformationen mit kontrollierten Transformations- und Prüfschritten verbindet. Zugleich muss empirisch untersucht werden, unter welchen Bedingungen ein solcher Ansatz technisch verarbeitbare und fachlich geeignete Ergebnisse erzeugt. Die vorliegende Arbeit entwickelt und evaluiert hierzu einen LLM-basierten Prototyp für die Transformation von Lebensmittel-Produktdaten. Die technische Verarbeitbarkeit wird für das umfassende Extraktionsergebnis betrachtet. Für das Referenzdatenfeld \textit{ProductSizing} beziehungsweise die Verpackungsgröße wird der LLM-basierte Ansatz zusätzlich einem regelbasierten Verfahren gegenübergestellt. Der regelbasierte Ansatz bildet dabei keine vollständige Alternative zur gesamten Produktextraktion, sondern eine feldspezifische Vergleichsbasis. Ergänzend werden Fehlerarten sowie Laufzeit, Tokenverbrauch und Verarbeitungskosten berücksichtigt. Die daraus abgeleiteten Ziele und Forschungsfragen werden im folgenden Abschnitt konkretisiert.

\subsection{Zielsetzung, Forschungsfragen und Abgrenzung}

Ziel der vorliegenden Arbeit ist die Konzeption, prototypische Implementierung und quantitative Evaluation eines Ansatzes zur automatisierten Transformation heterogener Lebensmittel-Produktdaten in ein bestehendes Zielschema. Hierzu wird mit dem CrawlDataTransformer ein IT-Artefakt entwickelt, das bereits gecrawlte Produktseiten vorverarbeitet, relevante Informationen mithilfe eines großen Sprachmodells extrahiert und die erzeugten Daten für die weitere Verarbeitung strukturiert, validiert und normalisiert. Das Vorgehen orientiert sich am Design-Science-Research-Paradigma, das die Entwicklung und Evaluation von IT-Artefakten zur Bearbeitung relevanter Problemstellungen in den Mittelpunkt stellt \cite{hevner2004design}.

Die Leistungsfähigkeit des LLM-basierten Ansatzes wird auf zwei voneinander zu unterscheidenden Ebenen betrachtet. Für das vollständige Produktextraktionsergebnis wird untersucht, in welchem Umfang die Modellausgaben syntaktisch verarbeitet, gegen das interne Produktschema validiert, normalisiert und durch die nachgelagerten Verarbeitungsschritte geführt werden können. Diese Untersuchung beschreibt die technische Verarbeitbarkeit der erzeugten Ergebnisse. Eine erfolgreiche Schema-Validierung belegt dabei nicht automatisch, dass sämtliche enthaltenen Produktinformationen fachlich korrekt oder vollständig sind.

Eine referenzbasierte fachliche Bewertung erfolgt für das Datenfeld \textit{ProductSizing}, das die Verpackungsgröße eines Produkts repräsentiert. Für dieses Feld werden die LLM-basierten Vorhersagen einer manuell erstellten Ground Truth gegenübergestellt. Zusätzlich wird ein regelbasierter Referenzansatz implementiert, der dieselben vorverarbeiteten Eingabedaten verwendet und ebenfalls eine Verpackungsgröße erzeugt. Dadurch können beide Verfahren unter gemeinsamen Referenz-, Normalisierungs- und Vergleichsbedingungen hinsichtlich ihrer Extraktionsqualität untersucht werden. Der regelbasierte Ansatz bildet dabei keine vollständige Alternative zur LLM-basierten Produktextraktion, sondern eine feldspezifische Baseline für den kontrollierten Vergleich.

Aus dieser Zielsetzung ergibt sich folgende übergeordnete Forschungsfrage:

\begin{quote}
Wie kann ein LLM-basierter Ansatz zur automatisierten Transformation heterogener Lebensmittel-Produktdaten in ein bestehendes Zielschema konzipiert und prototypisch umgesetzt werden, und welche technische sowie fachliche Leistungsfähigkeit zeigt dieser unter den untersuchten Bedingungen?
\end{quote}

Zur differenzierten Beantwortung wird die übergeordnete Forschungsfrage in drei Teilfragen untergliedert:

\begin{enumerate}
    \item In welchem Umfang erzeugt der LLM-basierte Ansatz syntaktisch verarbeitbare, schema-konforme und innerhalb der Transformationspipeline normalisierbare Produktextraktionsergebnisse?

    \item Welche fachliche Extraktionsqualität erreicht der LLM-basierte Ansatz für das Referenzdatenfeld \textit{ProductSizing}, und wie unterscheidet sie sich unter gemeinsamen Vergleichsbedingungen von der Qualität des regelbasierten Ansatzes?

    \item Welche Fehlerarten sowie Laufzeit-, Token- und Kostencharakteristika treten bei der prototypischen Verarbeitung auf, und welche Schlussfolgerungen ergeben sich daraus für die praktische Einsetzbarkeit des Ansatzes?
\end{enumerate}

Die erste Teilfrage wird anhand technischer Erfolgsstufen wie der Parsebarkeit, der Schema-Validierung, der Normalisierung und dem Abschluss der verpflichtenden Verarbeitungsschritte untersucht. Die zweite Teilfrage betrachtet die fachliche Übereinstimmung mit der Ground Truth und die Unterschiede zwischen beiden Extraktionsansätzen. Die dritte Teilfrage ergänzt die Qualitätsbewertung um beobachtete Fehlerarten und Effizienzaspekte. Die praktische Einsetzbarkeit wird dabei nicht aus einer einzelnen Kennzahl abgeleitet, sondern aus der gemeinsamen Betrachtung der technischen Verarbeitbarkeit, der fachlichen Extraktionsqualität und des erforderlichen Verarbeitungsaufwands.

Der Untersuchungsumfang ist auf die Transformation bereits vorhandener Crawl-Daten beschränkt. Die Entwicklung oder Veränderung des vorgelagerten Crawling-Systems ist ebenso wenig Bestandteil der Arbeit wie eine Erweiterung des vorgegebenen Zielschemas. Für die LLM-basierte Extraktion wird ein bestehendes externes Sprachmodell verwendet; das Training oder die domänenspezifische Feinabstimmung eines eigenen Modells wird nicht untersucht. Ebenso wird kein vollständiger regelbasierter Ersatz für die Produktextraktion entwickelt.

Da die manuell erstellte Ground Truth ausschließlich das Referenzdatenfeld \textit{ProductSizing} abdeckt, kann die fachliche Korrektheit und Vollständigkeit der übrigen Produktattribute nicht quantitativ gegen Referenzwerte bewertet werden. Für diese Attribute beschränkt sich die Untersuchung auf die technische Verarbeitbarkeit des gemeinsam erzeugten Produktextraktionsergebnisses. Die Evaluation bezieht sich zudem auf den verwendeten Datenbestand, die gewählte Promptkonfiguration, das eingesetzte Sprachmodell und die implementierten Regeln. Eine uneingeschränkte Übertragbarkeit der Ergebnisse auf andere Produktdomänen, Datenquellen, Sprachmodelle oder Konfigurationen wird daher nicht vorausgesetzt. Der entwickelte Prototyp dient der konzeptionellen und empirischen Untersuchung des Lösungsansatzes; Anforderungen an einen hochverfügbaren, horizontal skalierbaren Produktivbetrieb sind nicht Gegenstand der Arbeit.

\subsection{Aufbau der Arbeit}

Kapitel 3 schafft die theoretische Grundlage für die Untersuchung. Zunächst werden die für die Transformationsaufgabe relevanten Begriffe und Zusammenhänge aus den Bereichen Datenintegration, Datenqualität und Informationsextraktion eingeführt. Anschließend werden regelbasierte und LLM-basierte Extraktionsverfahren sowie geeignete Grundlagen ihrer Evaluation betrachtet. Die Einordnung verwandter Arbeiten führt die dargestellten Ansätze zusammen und leitet daraus die durch die vorliegende Arbeit adressierte Forschungslücke ab.

Auf dieser Grundlage entwickelt Kapitel 4 den konzeptionellen Lösungsansatz. Aus der Analyse der Quelldaten, des vorgegebenen Zielschemas und des bestehenden Systemkontexts werden die funktionalen und nicht-funktionalen Anforderungen an den Prototyp abgeleitet. Darauf aufbauend werden die Architekturziele, die zentralen Entwurfsentscheidungen sowie die für Verarbeitung, Persistenz, Integration, Evaluation und Fehlerbehandlung maßgeblichen Architekturkonzepte beschrieben. Kapitel 5 stellt anschließend die prototypische Umsetzung der für die Forschungsfragen relevanten Mechanismen dar. Im Mittelpunkt stehen die mehrphasige Transformationsverarbeitung, die gemeinsame HTML-Vorverarbeitung, die LLM-basierte Extraktion, der regelbasierte Referenzansatz für \textit{ProductSizing} sowie die nachgelagerte Validierung, Normalisierung und Evaluationsunterstützung.

Kapitel 6 untersucht das entwickelte Artefakt anhand eines manuell annotierten Testdatensatzes. Nach der Darstellung des Evaluationsdesigns werden zunächst die technische Verarbeitbarkeit und die Effizienz der vollständigen LLM-basierten Produktextraktion bewertet. Anschließend erfolgt für das Referenzdatenfeld \textit{ProductSizing} der fachliche Vergleich mit dem regelbasierten Ansatz. Die Ergebnisse werden abschließend hinsichtlich ihrer praktischen Einsetzbarkeit, ihrer Aussagekraft und ihrer Grenzen diskutiert. Kapitel 7 führt die wesentlichen Erkenntnisse zusammen, beantwortet die Forschungsfragen und leitet aus der kritischen Würdigung mögliche weiterführende Arbeiten ab. Der Anhang ergänzt die Darstellung um gekürzte Beispiele der verwendeten Eingabe- und Zielartefakte.